\section*{Limitations}
We aim to provide broad coverage of methods that use verbal feedback within the agent's lifecycle. Within each category, we focus on representative papers rather than attempting an exhaustive enumeration. Additionally, while our taxonomy offers a unifying framework, it may oversimplify methods that span multiple roles of language. Several approaches do not fit neatly into a single branch, since verbal feedback can simultaneously serve as a grounding signal, a deliberative mechanism, and a learning signal. Our taxonomy should therefore be understood as an organizing lens rather than a strict partition of the literature. Finally, our survey focuses primarily on methods in which verbal feedback is explicit and central to the agent's process, and therefore does not fully cover adjacent work where language plays a more auxiliary role. \added{Language-based feedback itself carries failure modes that this survey can name but not yet bound: reward-specification failures when language compiles to the wrong objective, unreliable interpretation of feedback by the consuming model, weak grounding between verbal descriptions and environment state, and the transfer of language-guided policies to real-world settings. Societal impact follows the same channels: when feedback is incorrect, biased, adversarial, or poorly grounded, harm falls on users of high-stakes deployments such as coding assistants, robotics, clinical decision support, and education, and on the people affected by those systems' outputs. Sections~\ref{sec:adversarial_verbal_feedback} and~\ref{sec:theoritical_foundation} discuss mitigation directions.} AI writing tools were used for language editing and proofreading only; all research contributions are the authors' own.
% R2 will expand this section (in red) with: reviewer-named failure modes of language-based
% feedback (reward-specification failures, unreliable language interpretation, weak grounding,
% real-world transfer) and who-is-harmed societal impact for high-stakes deployments.
